\documentclass{article}
\usepackage[utf8]{inputenc}

\usepackage{amsmath, amssymb, amsthm, amsfonts}
\usepackage{subcaption, enumerate}
\usepackage{bbold}
\usepackage{polynom}
\usepackage{xfrac}
\usepackage{graphicx}
\usepackage{mathtools}
\usepackage{enumitem}

\usepackage{tikz}

\theoremstyle{definition}
\newtheorem{definition}{Definition}
\newtheorem{example}{Example}[definition]
\newcommand{\dab}
{
	\begin{tikzpicture}[scale=0.04]
		% head
		\draw[thick] (0, 0) circle (3cm);

		% bent arm
		\draw[thick, rotate around={15 : (0, 0)}] (-3, 0) -- (-3.7, -2) -- (4.1, -1.1);

		% extended arm
		\draw[thick, rotate=20] (3, 0) -- (7, 2);
	\end{tikzpicture}
}
\newcommand{\st}{\text{ s.t. }}

\makeatletter
\renewcommand*\env@matrix[1][*\c@MaxMatrixCols c]{%
  \hskip -\arraycolsep
  \let\@ifnextchar\new@ifnextchar
  \array{#1}}
\makeatother

\usepackage{fullpage}

\title{Linear Transforms Day 7 \\ with Dr Nevard}
\author{Notes by Laura Gallo}
\date{October 6, 2020}

\begin{document}
\maketitle

\section{Finding Inverses of Matricese}
We can use Gaussian elimination to find inverses of matrices. For example:
\begin{equation*}
	\left(
		\begin{matrix}[cc|cc]
			1 & 2 & 1 & 0 \\
			3 & 4 & 0 & 1
		\end{matrix}
	\right)
\end{equation*}
\begin{equation*}
	\left(
		\begin{matrix}[cc|cc]
			1 & 2 & 1 & 0 \\
			0 & -2 & -3 & 1
		\end{matrix}
	\right)
\end{equation*}
\begin{equation*}
	\left(
		\begin{matrix}[cc|cc]
			1 & 0 & -2 & 1 \\
			0 & -2 & -3 & 1
		\end{matrix}
	\right)
\end{equation*}
\begin{equation*}
	\left(
		\begin{matrix}[cc|cc]
			1 & 0 & -2 & 1 \\
			0 & 1 & -\frac{3}{2} & \frac{1}{2}
		\end{matrix}
	\right)
\end{equation*}
The right side of the augmented matrix is the inverse of the original matrix; ie,
\begin{equation*}
	\left(
		\begin{matrix}
			1 & 2 \\
			3 & 4
		\end{matrix}
	\right)
	\left(
		\begin{matrix}
			-2 & 1 \\
			-\frac{3}{2} & \frac{1}{2}
		\end{matrix}
	\right) = \left(
		\begin{matrix}
			1 & 0 \\
			0 & 1
		\end{matrix}
	\right)
\end{equation*}
Note that this is an inefficient way to solve systems of equations.
\end{document}
